% Fragmento público generado desde propietarios íntegros.
% Residencia: 52.2.2.
% Fuente: ../incoming/radion_monografia_20260827/manuscrito/sections/03_cierre_exacto.tex:51-103

\paragraph{Densidad orientada \(20/81\) del sector activo \(90/120\)}

El selector trítico de fase separa
\[
H_+=\{1,4,7\},\qquad H_-=\{2,5,8\},\qquad H_0=\{3,6,9\}.
\]
Hay seis fases activas y quince parejas no ordenadas:
\[
H_{\mathrm{act}}=H_+\cup H_-,\qquad
\left|\binom{H_{\mathrm{act}}}2\right|=\binom62=15.
\]
El modo activo y su extensión a ocho posiciones son
\begin{equation}
N_6=6\binom62=90,\qquad
N_8=8\binom62=120.
\label{eq:90-120-incidencia}
\end{equation}
Dos hojas conjugadas del sector activo, normalizadas por la fibra ambiente,
dan
\begin{equation}
\rho_{\mathrm{tw}}
=\frac{2N_6}{|\F_3^6|}
=\frac{2\cdot90}{729}
=\frac{180}{729}
=\boxed{\frac{20}{81}}.
\label{eq:rho-twoway}
\end{equation}

\begin{proposition}[Unicidad relativa de la densidad]
Fijados el sector activo \(N_6=90\), las dos hojas conjugadas y la
normalización por \(|\F_3^6|=729\), la densidad orientada es necesariamente
\(20/81\).
\end{proposition}
\begin{proof}
La cardinalidad transportada es \(2N_6=180\). La normalización por el ambiente
da \(180/729\); dividir numerador y denominador por \(9\) produce \(20/81\).
No interviene ninguna comparación con \(180/\pi\) ni con \(\epsR\).
\end{proof}

La lectura histórica como proyección sobre una red \(1/81\) es un control de
rigidez útil, pero ya no es el fundamento: el coeficiente se obtiene antes
del cierre. La ablación de una hoja produce \(10/81\), y la salida cambia en
orden \(10^{-5}\); por tanto, el factor dos no es decorativo.

La sección torsional es
\begin{equation}
\Phi_T=\rho_{\mathrm{tw}}\Delta_4,
\qquad
S_T=1+\Phi_T.
\label{eq:seccion-torsional}
\end{equation}

